% !TeX root = ../../Book3.tex
% 纲要标题：### 14.3 切空间与 Jacobian 矩阵
% 纲要位置：工作记录/计划纲要.md，第 2417 行。
% 纲要细目（写作范围）：
% * 极大理想模平方与余切空间
% * 仿射方程的切空间计算
% * 一般素点的剩余域微分方向；不可分闭点与基变换后的幂零结构

\section{切空间与 Jacobian 矩阵}
\label{b3:sec:1403}

一条曲线在一个点附近可能有几条分支，也可能含有重根；只看点集不能分辨这些现象。把点的坐标改成 \(a_i+v_i\varepsilon\)、其中 \(\varepsilon^2=0\)，方程留下的一次关系便限制了允许的方向。先从有理点说明这个计算，再讨论剩余域不是基域的情形。


\subsection{极大理想模平方与余切空间}
\label{b3:subsec:1403:01}

\begin{definition}\label{b3:def:zariski-tangent-space}
设 \(k\) 为域，\(B\) 为有限型交换 \(k\) 代数，\(a:B\to k\) 为一个 \(k\) 代数同态，\(\mathfrak m=\ker a\)。以 \(a\) 赋予 \(k\) 一个 \(B\) 模结构。向量空间
\[
T_a(B)=\operatorname{Der}_k(B,k)
\]
称为 \(B\) 在这个 \(k\) 有理点处的\term[zariski-qiekongjian]{Zariski 切空间}[Zariski tangent space]；\(\mathfrak m/\mathfrak m^2\) 是该点的余切空间。这里的 \(B\) 可以含有非零幂零元，不要求它是某个点集的约化坐标环。
\symidx{\(T_a(B)\)：有理点处的 Zariski 切空间}
\end{definition}

\begin{proposition}\label{b3:prop:tangent-cotangent-rational}
设 \(k\) 为域，\(B\) 为交换 \(k\) 代数，\(a:B\to k\) 为 \(k\) 代数同态，\(\mathfrak m=\ker a\)。则有自然同构
\[
\mathfrak m/\mathfrak m^2\xrightarrow{\sim}
k\otimes_B\Omega_{B/k},\qquad
\bar b\longmapsto1\otimes \mathrm{d}b,
\]
以及
\[
\operatorname{Der}_k(B,k)
\cong\operatorname{Hom}_k(\mathfrak m/\mathfrak m^2,k).
\]
若以局部环 \(B_{\mathfrak m}\) 代替 \(B\)，所得余切空间及切空间相同。
\end{proposition}
\begin{proof}
对 \(b,c\in\mathfrak m\)，\(1\otimes \mathrm{d}(bc)=0\)，故第一个映射良定义。每个 \(b\in B\) 都写成 \(a(b)+\bigl(b-a(b)\bigr)\)，其中后项属于 \(\mathfrak m\)，所以这个映射满射。反向定义
\[
D:B\longrightarrow\mathfrak m/\mathfrak m^2,
\qquad D(b)=b-a(b)\pmod{\mathfrak m^2}.
\]
将 \(b=a(b)+b_0\)、\(c=a(c)+c_0\) 相乘，模掉 \(\mathfrak m^2\) 后得到
\(D(bc)=a(b)\cdot D(c)+a(c)\cdot D(b)\)。故 \(D\) 是取值于剩余域模的 \(k\) 导子，由微分模泛性质诱导逆映射。第二个同构由同一泛性质及第一个同构得到。

若 \(s\notin\mathfrak m\)，它在 \(\mathfrak m/\mathfrak m^2\) 上的乘法就是非零标量 \(a(s)\)，已经可逆。因此局部化不改变该商；由局部化正合性，它又等于
\(\mathfrak mB_{\mathfrak m}/(\mathfrak mB_{\mathfrak m})^2\)。导子的分式延拓公式也给出相同的切空间。
\end{proof}

例如，在 \(k[x_1,\ldots,x_n]_{(x_1,\ldots,x_n)}\) 的原点，记极大理想为 \(\mathfrak m\)。消失于原点的函数模 \(\mathfrak m^2\) 后，只留下 \(\sum_i c_ix_i\) 这样的一次部分；\(\mathfrak m/\mathfrak m^2\) 保存的是函数的一阶线性信息。切向量 \(v=(v_i)\) 则是作用在这些一次部分上的线性泛函，将 \(\sum_i c_ix_i\) 送到 \(\sum_i c_iv_i\)。这就是切空间与余切空间互为对偶的具体含义。

上述导子还有直接的点的解释。给定 \(D\in T_a(B)\)，映射
\[
B\longrightarrow k[\varepsilon]/(\varepsilon^2),\qquad
b\longmapsto a(b)+D(b)\cdot\varepsilon
\]
是一个模 \(\varepsilon\) 后等于 \(a\) 的 \(k\) 代数同态；乘法相容性恰是导子的公式。反过来，每个这样的同态都唯一给出一个 \(D\)。因此，同一个切向量既可视为导子，也可视为余切空间上的线性泛函，或点向双数环中的一阶延伸。

\subsection{仿射方程的切空间计算}
\label{b3:subsec:1403:02}

\begin{theorem}\label{b3:thm:jacobian-tangent-space}
设 \(k\) 为域，\(B=k[X_1,\ldots,X_n]/(f_1,\ldots,f_s)\)，且 \(a=(a_1,\ldots,a_n)\in k^n\) 满足所有 \(f_i(a)=0\)。以 \(X_j\mapsto a_j\) 的求值同态 \(B\to k\) 定义该点的切空间 \(T_a(B)=\operatorname{Der}_k(B,k)\)。令
\[
J_f(a)=\left(\dfrac{\partial f_i}{\partial X_j}(a)\right)_{
1\le i\le s,\,1\le j\le n}.
\]
则在变量坐标给出的识别下，
\[
T_a(B)=\ker\bigl(J_f(a):k^n\to k^s\bigr),
\qquad \dim_kT_a(B)=n-\operatorname{rank}J_f(a).
\]
这个矩阵称为所给方程组的\term[jacobian-juzhen]{Jacobian 矩阵}[Jacobian matrix]，亦称雅可比矩阵\termidx[yakebi-juzhen]{雅可比矩阵}。\footnote{“雅可比矩阵”见菲赫金哥尔茨《微积分学教程》第一卷\cite[第六章“函数矩阵(雅可比矩阵)的乘法”]{Fichtenholz2005CalculusChineseI}。}
\end{theorem}
\begin{proof}
由双数值点的描述，一个方向 \(v=(v_j)\) 对应 \(X_j\mapsto a_j+v_j\varepsilon\)。它能从多项式环下降到 \(B\)，当且仅当每个 \(f_i\) 的像为零。因为 \(\varepsilon^2=0\)，二次以上的改变量全部消失，故
\[
f_i(a_1+v_1\varepsilon,\ldots,a_n+v_n\varepsilon)
=f_i(a)+\varepsilon\sum_j
\dfrac{\partial f_i}{\partial X_j}(a)v_j.
\]
由 \(f_i(a)=0\)，这一条件恰为 \(J_f(a)v=0\)，得到核的描述及维数公式。

\cref{b3:cor:differential-presentation}还将这一计算与环本身的微分模联系起来：展示 \(B^s\to B^n\to\Omega_{B/k}\to0\) 的关系矩阵为 \(J_f^{\mathsf t}\)。张量到 \(k\) 后仍右正合，再取 \(k\) 线性对偶，得到
\[
0\longrightarrow\operatorname{Hom}_k(k\otimes_B\Omega_{B/k},k)
\longrightarrow k^n\xrightarrow{J_f(a)}k^s.
\]
前一命题把左端自然识别为切空间。因此坐标与方程展示虽会改变矩阵，所得导子空间始终由 \(B\) 及这个有理点决定。
\end{proof}

例如设 \(k\) 为域，在 \(B=k[x,y]/(xy-1)\) 中，每个有理点 \((a,b)\) 都有 \(ab=1\)，Jacobian 行为 \((b,a)\)，秩为一，所以切空间由 \(bv_x+av_y=0\) 给出，是一维的。在尖点曲线 \(y^2=x^3\) 的原点，Jacobian 行 \((-3x^2,2y)\) 为零，故切空间是整个 \(k^2\)，任意特征都如此。曲线 \(xy=0\) 的原点也有二维切空间，虽然它有两条分支，而尖点只有一条。这说明切空间维数能发现某些异常，却不能单独数出分支。

\ProseParagraphBreak
尖点的参数化 \(t\mapsto(t^2,t^3)\) 把源直线原点的一阶点 \(t=v\varepsilon\) 送到 \(\bigl((v\varepsilon)^2,(v\varepsilon)^3\bigr)=(0,0)\)。因此它在原点诱导的切映射为零，源点的一阶变化在像中全部消失。

竖直方向则显示 Zariski 一阶切方向与切锥约化方向之间的差别。映射 \(x\mapsto0,y\mapsto\varepsilon\) 在 \(k[\varepsilon]/(\varepsilon^2)\) 中满足尖点方程，因此给出一个切向量。若提升到 \(k[\varepsilon]/(\varepsilon^3)\)，保持这一阶值就必须写成 \(x=a\varepsilon^2\)、\(y=\varepsilon+b\varepsilon^2\)。此时 \(x^3=0\)，而 \(y^2=\varepsilon^2\ne0\)，所以没有这样的提升。Zariski 切空间只检查一次关系；切锥 \(Y^2=0\) 则已保留最低二次关系，其约化支集只有水平直线。因而一阶允许的方向未必能延续成更高阶曲线。

\ProseParagraphBreak
还须保留方程中的幂零结构。\(k[t]/(t^2)\) 只有一个域值点，但该点切空间为一维；约化后的坐标环 \(k\) 则有零切空间。两者的点集相同，一阶延伸却不同。若已经把理想替换成根理想再计算，便无法恢复丢掉的方向。

\ProseParagraphBreak
Jacobian 计算还可以寻找参数族中切空间异常增大的位置。以 \(t\) 为参数，在特征零域 \(k\) 上考虑曲线族
\[
F_t(x,y)=y^2-x^3-t=0.
\]
固定参数后，切空间维数为二恰好要求 \(F_t=0\)、\(\partial F_t/\partial x=-3x^2=0\)、\(\partial F_t/\partial y=2y=0\)。在 \(k[t,x,y]\) 中，相应理想为
\[
(y^2-x^3-t,-3x^2,2y)=(y,x^2,t).
\]
右边的包含关系可直接由 \(t=y^2-x^3-F_t\) 核验，反向也由各生成元的表达式得到。这个 Jacobian 理想记录异常位置的代数结构，其商中还保留 \(x^2=0\) 的加厚信息；若只求底层异常点集，则取根理想得到 \((x,y,t)\)，只剩原点。要判断哪些参数会产生异常，还须消去 \(x,y\)，收缩到参数环得到
\[
(y,x^2,t)\cap k[t]=(t).
\]
这是因为 \(t\) 已属于该理想，而属于它的参数多项式必在原点取值为零。于是异常只在 \(t=0\) 出现；其余参数的每个几何点都有一维切空间。

\subsection{非完美基域上的切空间}
\label{b3:subsec:1403:03}

这里先比较一般素点处的两个一阶对象，差别并不限于不完美基域。设 \(k\) 为域，\(B\) 为交换 \(k\) 代数，\(\mathfrak p\in\operatorname{Spec}B\)，记 \(\kappa=\kappa(\mathfrak p)\)。剩余域 \(\kappa\) 未必等于 \(k\)。局部环自身的余切空间
\(\mathfrak pB_{\mathfrak p}/(\mathfrak pB_{\mathfrak p})^2\)
与相对微分的纤维 \(\kappa\otimes_B\Omega_{B/k}\) 出自不同构造。对 \(B_{\mathfrak p}\to\kappa\) 应用余法正合列及微分局部化，得到下列 \(\kappa\) 向量空间的右正合列：
\begin{equation}\label{b3:eq:cotangent-residue-sequence}
\begin{tikzcd}[column sep=small,row sep=0.8em]
\dfrac{\mathfrak pB_{\mathfrak p}}{(\mathfrak pB_{\mathfrak p})^2}
  \arrow[r]
&\kappa\otimes_B\Omega_{B/k}
  \arrow[r]
&\Omega_{\kappa/k}\arrow[r]&0\\
{\begin{gathered}\text{局部环自身}\\\text{的一阶结构}\end{gathered}}
&{\begin{gathered}\text{相对于基域 }k\\\text{的一阶结构}\end{gathered}}
&{\begin{gathered}\text{剩余域扩张}\\\text{的一阶结构}\end{gathered}}&
\end{tikzcd}
\end{equation}
在 \(k\) 有理点处，\cref{b3:prop:tangent-cotangent-rational}证明左箭头为同构。一般点处的左箭头不必单射，右端也可能非零；它记录了剩余域扩张本身贡献的微分。记左箭头的像为 \(W\)，则中间空间包含子空间 \(W\)，商空间为 \(\Omega_{\kappa/k}\)。向量空间的短正合列可以选基分裂，所以中间空间可非典范地写成 \(W\oplus\Omega_{\kappa/k}\)；但没有自然指定的补空间，而且 \(W\) 只是左端的一个商，不能直接把整个局部余切空间与 \(\Omega_{\kappa/k}\) 作直和。

即使 \(k\) 完美，在 \(B=k[t]\) 的一般点 \(\mathfrak p=(0)\) 处也有 \(\kappa=k(t)\)；由多项式微分与局部化公式，\(\Omega_{\kappa/k}=k(t)\,\mathrm{d}t\ne0\)。此时局部环是域，局部余切空间为零，相对微分纤维却仍有一维。非完美基域上的额外现象是：连闭点的有限剩余域扩张也可能贡献这样的非零微分。

\begin{example}\label{b3:exm:inseparable-tangent}
设 \(p\) 为素数、\(k=\mathbb F_p(s)\)，其中 \(s\) 为超越元。令
\[
L=k[u]/(u^p-s).
\]
由于 \(s\) 在 \(k\) 中不是 \(p\) 次幂，\(u^p-s\) 不可约，故 \(L\) 为域。其唯一局部环的极大理想为零，局部环余切空间也为零。但微分展示给出
\[
\Omega_{L/k}=L\,\mathrm{d}u,
\]
因为 \(s\) 是基域中的常数，且 \(k\) 的特征为 \(p\)，故 \(\mathrm{d}(u^p-s)=0\)。该点的剩余域就是 \(L\)，所以相对微分纤维 \(L\otimes_L\Omega_{L/k}=\Omega_{L/k}\) 非零，完全来自剩余域扩张的不可分性，并非来自局部极大理想。若扩域到 \(K=k(s^{1/p})\)，则
\[
K\otimes_kL\cong K[v]/(v^p),\qquad v=u-s^{1/p},
\]
其唯一 \(K\) 有理点由 \(v\mapsto0\) 给出，局部余切空间 \((v)/(v^2)\) 以 \(v\) 的类为基，故该点的切空间为一维。这里 \(v\) 的像是非零幂零元：原来的环 \(L\) 是域，改变基域后却出现了这个加厚点。
\end{example}

不可约性的一个直接检验如下。在代数闭包中，\(T^p-s=(T-\alpha)^p\)。若 \(\alpha\) 的最小多项式次数 \(r<p\)，它只能为 \((T-\alpha)^r\)，其 \(T^{r-1}\) 系数使 \(r\alpha\in k\)，从而 \(\alpha\in k\)，与 \(s\) 非 \(p\) 次幂矛盾。最后一事实可由有理函数在 \(s=0\) 处的零点阶数核验：\(p\) 次幂的阶数是 \(p\) 的倍数，\(s\) 的阶数则为一。

\ProseParagraphBreak
这一区别在不完美基域上尤其必要。局部环是域，因而在环论意义下最简单，却不能据此断言它相对于给定基域光滑。后面的正则局部环与光滑性判据将分别处理环自身的性质和扩张基域后的性质；本章的微分计算已经显示两者为何不能合并。
